\documentclass{article}
\usepackage[T1]{fontenc}
\usepackage{lmodern}
\usepackage{graphicx}
\usepackage{amsmath,amssymb}
\usepackage[a4paper,margin=2cm]{geometry}
\usepackage[absolute,overlay]{textpos}
\setlength{\TPHorizModule}{1mm}
\setlength{\TPVertModule}{1mm}
\usepackage[indonesian]{babel}
\usepackage{hyperref}
\setlength{\emergencystretch}{2em}
% unit: Halaman 12

\begin{document}

% === Halaman 12 (python/native) ===

Algoritma Knapsack Kunci-Publik

• Algoritma superincreasing knapsack adalah algoritma yang lemah, karena 

cipherteks dapat didekripsi menjadi plainteksnya secara mudah dalam 

waktu lanjar. 

• Algoritma non-superincreasing knapsack atau normal knapsack adalah 

kelompok algoritma knapsack yang sulit (dari segi komputasi) karena 

membutuhkan waktu dalam orde eksponensial untuk memecahkannya.

• Namun, superincreasing knapsack dapat dimodifikasi menjadi non-

superincreasing knapsack dengan menggunakan kunci publik (untuk 

enkripsi) dan kunci privat (untuk dekripsi). 

12

\end{document}
